%%==============================================================================
%% Estilo de la hoja | @file /informe_latex/format.tex
%%==============================================================================

% Creamos una variable y guardamos de forma limpia el parskip actual (los 10ex)
\newlength{\parskipGlobal}
\setlength{\parskipGlobal}{1ex}


% Adjust spacing after the chapter title.
\titlespacing*{\chapter}{0cm}{-2.0cm}{0.50cm}
\titlespacing*{\section}{0cm}{0.50cm}{0.25cm}

% Indent 
\setlength{\parindent}{0pt}
\setlength{\parskip}{\parskipGlobal}


\newtcolorbox{problem}{
	enhanced,
	colback = black!5,
	coltitle = black,
	boxrule = 0pt,
	frame hidden,
	borderline west = {0.5mm}{0.0mm}{black},
	fontupper = \normalfont,
	title = {},
	breakable,
	before skip = 3ex,
	after skip = 3ex,
	%
	% Lee dinámicamente el parskip global actual (1ex) y lo fija al inicio del texto interno
	before upper = {\setlength{\parskip}{\parskipGlobal}}
}
\tcbuselibrary{skins, breakable}

% Ajuste global para listas (acorta el tabulador horizontal).
\setlist[itemize]{
	leftmargin = 2.5em,  % Distancia desde el margen izquierdo general hasta el texto
	labelsep = 0.5em,    % Distancia entre el puntito/cuadradito y el texto
	itemindent = 0pt     % Alineación base del ítem
}
